\section{Positive-source inputs and the topology of conditioning}
\label{sec:v92-inputs}
The following statements are used in precisely the forms proved in the
retained source. They are not inferred from source preservation or
numerical diagnostics. Write $\mathcal C(M)=\mathcal C_{m,R,n,k}(M)$,
$P=m^2p$, $G=\mathcal L_{m,R}$ and $Q=32/31$. For $q>1$ set
\[
 \mathcal N_q(u;J)=\sup_{E\subset J,\,0<|E|<\infty}
                         |E|^{-1+1/q}\int_E|u|.
\]
This is the measurable-set norm of weak $L^q$, not strong $L^q$.
The finite exponent $Q$ is not the pointwise endpoint.

\begin{inputstatement}[Positive scalar remainder]
\label{input:v92-scalar}
For each sufficiently large fixed $B$,
\begin{equation}\label{eq:v92-positive-input}
 P-G=b_B+e_B,\qquad b_B=m^2b^{\varepsilon(B)}\ge0,\qquad
 \varepsilon(B)=A_0B^{-1/12},
\end{equation}
where
\[
 \limsup_m\sup_{R,n,k}\esssup_{\mathcal C(M)}|e_B|
                                       \le C_MB^{-1/192}.
\]
The original finite-peak reference has $|G|\le D_M$ on the central
class. The constants are uniform in radius and exact labels. All
auxiliary bands are fixed before the collision limit and enlarged
only afterwards, with $B$ outermost.
\end{inputstatement}
This is Supplement Theorem~\ref{supp-thm:v51-positive-error}, in the
normalization restated in Supplement equation
\eqref{supp-v91archive:eq:v91-positive-remainder}.

\begin{inputstatement}[The unweakened critical limit and original union]
\label{input:v92-critical}
For each fixed $M,L>0$,
\[
 \sup_{R,n,k}\sup_{|J|=L}
      \mathcal N_Q((P-G)\one_{\mathcal C(M)};J)\longrightarrow0.
\]
The same critical convergence holds after
$t=m\bar\tau_R+\sqrt m\,u$ on $u\in[-M,M]$, with the central
indicator retained. Thus the full-window measure is $dt/\sqrt m$.
For $z=|\nabla_qF|$ and total maximal regular residence $\ell$, define
\[
 U_{a,\delta}(t)=\int\one_{\{z<a\ \mathrm{or}\ \ell<\delta\}}\,d\mu_t.
\]
At the cofinal cutoffs $a=h^{1/8}$, $\delta=h^{1/4}$,
\begin{equation}\label{eq:v92-union-input}
 \limsup_m\sup_{R,n,k}\sup_{|J|=L}
 \mathcal N_Q(m^2U_{a,\delta}\one_{\mathcal C(M)};J)
                                   \le C_{M,L}h^{31/4096}.
\end{equation}
Every original weighted selector with smaller cutoffs is dominated by
this same nonnegative upper source, since $0\le W^\varepsilon\le1$.
The definition and normalization of the weighted target are unchanged.
\end{inputstatement}
These are Supplement Theorem~\ref{supp-v91archive:thm:v91-raw} and
Theorem~\ref{supp-v91archive:thm:v91-exceptional}. Their proofs use
the complete first-exit surface, including every positive residence
scale, before taking the critical norm. The inherited finite-norm
conclusions do not assume the original essential-height charge is zero.

\begin{inputstatement}[The positive path remainder on common roofs]
\label{input:v92-path}
Let $\mathcal X=C([0,1],\R^4)$. The actual collision path is tied to
its own terminal record, with exactly the convention of Supplement
Section~\ref{supp-sec:v52-path-inversion}. Its common conditional
kernel is $\mathsf Q_t$, and its centered Gaussian bridge reference
is $\mathsf W_R$. With $\mathbf P=P\mathsf Q_t$, the original guard
restriction has a positive path kernel $0\le\mathbf b_B\le\mathbf P$
with mass $b_B$. The same finite path measures satisfy
\begin{equation}\label{eq:v92-path-input}
 \mathbf P-G\mathsf W_R=\mathbf b_B+\mathbf E_B,\qquad
 \limsup_m\sup_{R,n,k}\esssup_{\mathcal C(M)}
             \|\mathbf E_B\|_{\mathrm{BL}^*}\le C_MB^{-1/192}.
\end{equation}
The bounded-Lipschitz class has supremum norm and Lipschitz constant
at most one and includes the constant one. Thus the norm of a positive
finite measure is its mass. The scalar remainder may be taken as
$e_B=\mathbf E_B(1)$ when both inputs are used.
\end{inputstatement}
This is Supplement Theorem~\ref{supp-thm:v52-path-remainder}, with
its cylinder inversion, variation tightness and endpoint-budget
hypotheses retained. Supplement equation
\eqref{supp-eq:v73-pointwise-norm-comparison} records the elementary
comparison subsequently used here. The path input is additional to
the scalar input; it is not deduced from scalar weak convergence.

For each fixed $m,R,n,k$, the original Borel coarea construction fixes
common roof kernels outside a null set. A countable path norming class
makes the norms in \eqref{eq:v92-path-input} measurable on the same
version. Parameter uniformity always means $\sup_R\esssup_t$; no
single null set for all radii is asserted. Conditional-kernel results
below are restricted to $G\ge d>0$, and moving-set results additionally
require a fixed positive reference-mass floor. No statement replaces
these quantified restrictions by mere positivity or assigns a
normalized conditional law to an arithmetic zero class.
